\section{Pairing and the gap along the clock}
\label{sec:gap}

\subsection{The pair channels}

A pair of fermions is created by an operator $\psi^{T}X\psi$ with a $16\times16$ matrix $X$; the
choice of $X$ is the \emph{channel} of the pair. The earlier paper classified the channels and
computed, for each, the \emph{exchange amplitudes} $F_S$ and $F_P$ of the radial mode and of the
phase mode: the matrix elements of the exchange of one boson between two fermions of opposite
momenta in the internal state of the channel. It found that the radial mode attracts in every
channel, and that the phase mode attracts in the four \emph{vector pairs along the extra axes},
\begin{equation}
\psi^{T}\Cm\gamma^{a}\psi,\qquad a\in\{0,5,6,7\},
\label{eq:pairs}
\end{equation}
and repels in their partners that are built with two Dirac matrices. Here
$\Cm=\SIG{\gamma^{0}\gamma^{1}\gamma^{2}\gamma^{3}}{\gamma^{4}\gamma^{0}\gamma^{1}\cdots\gamma^{7}}$
is the matrix that turns each Dirac matrix into minus its transpose\SIG{ (Sec.~\ref{sec:dirac})}{}, and the index $a$ labels the
four extra axes. The pair along the hidden axis, $a=0$, has spin 0 in ordinary spacetime; it is the
\emph{hidden-axis singlet} of hypothesis H1.

The amplitudes are computed from the Dirac matrices of the frame, which the deflation does not
change, and each is proportional to the square of the coupling. For fields that do not depend on
$x5,x6,x7$ the amplitude of a channel at the time $x4$ is therefore
\begin{equation}
\begin{aligned}
F_{S,P}(x4)={}&\gVS(x4)^{2}\\
&\times\big(F_{S,P}\text{ of the flat spacetime at }\gVS=1\big).
\end{aligned}
\label{eq:amplitudes}
\end{equation}
The program computes the amplitudes of the 15 representative channels of the earlier paper again
from the matrices of the deflating spacetime, at three momenta and 13 scattering angles; the numbers
agree with those of the earlier paper to every digit recorded there,
and the amplitudes of the hidden-axis channel reproduce its closed forms to $1.1\times10^{-15}$. The
conclusion of the earlier paper about the channels therefore stands at every moment
(Fig.~\ref{fig:kernels}): the most attractive channel is the vector pair along an extra axis.

\begin{figure*}
\includegraphics[width=0.84\textwidth]{fig_kernels}
\caption{\label{fig:kernels}Exchange amplitudes in three pair channels, divided by the square of the
coupling of the moment, against the scattering angle, for momenta of the magnitude $m$ in the plane
of $x1$ and $x2$. Lines: deflating $7+1$ spacetime; open symbols: deflating $4+4$ spacetime. (a) The
radial mode attracts in all three channels. (b) The phase mode attracts in the vector pair along the
hidden axis and repels in its partner built with two Dirac matrices. The amplitudes are those of flat
spacetime: the deflation changes their overall strength and nothing else.}
\end{figure*}

\subsection{The gap grows much faster than the coupling}

In matter of a given density the attraction opens a gap $\Delta_F$ at the Fermi momentum $\kF$: the
energy that it costs to break a pair. The gap equation of the earlier paper, which generalizes that
of Bardeen, Cooper and Schrieffer \cite{BCS1957}, holds at each moment with the couplings of that
moment. Its solution as a function of the clock is the \emph{adiabatic history} of the paired state:
the state that matter would have if it were always in equilibrium. (Section~\ref{sec:quench}
examines when it can be.) The strength of the pairing is measured by the dimensionless number
$\lambda=N(0)\,W(\kF,\kF)$, the density of states at the Fermi surface times the kernel of the gap
equation there; in the vector channel it is the sum $\lambda_S+\lambda_P$ of the contributions of the
two modes. Since the kernel is proportional to $\gVS^{2}$,
\begin{equation}
\lambda\ \propto\ e^{\td}.
\end{equation}

The program solves the gap equation at 34 moments for each of four sets of the mass $m_s$ and the
Fermi momentum (Table~\ref{tab:gap}, Fig.~\ref{fig:gap}). There are three results.

\begin{table*}
\caption{\label{tab:gap}The gap along the deflation clock for two of the four sets computed. $\gVS$
and $\lambda$: coupling and pairing strength of the moment. $\EF$: Fermi kinetic energy. $T_c$:
critical temperature. Rate of growth: $d\ln\Delta_F/d\td$; the last column, its product with
$\lambda$, is given where the gap is between $10^{-6}$ and $0.02$ of $\EF$.}
\begin{ruledtabular}
\begin{tabular}{cccccccc}
$m_s/m$, $\kF/m$ & $\td$ & $\gVS$ & $\lambda$ & $\Delta_F/\EF$ & $T_c/\Delta_F$ & rate of growth & $\lambda\,\times$ rate \\
\hline
{{ROWS:tbl_gap}}
\end{tabular}
\end{ruledtabular}
\end{table*}

\emph{The order of the channels does not change.} At every moment the gap of the vector channel is
at least as large as the gap that the radial mode alone would give, and that at least as large as
the gap of the channel in which the phase mode repels. The gap never decreases.

\emph{The critical temperature follows the gap.} Wherever the gap is below 5 per cent of the Fermi
kinetic energy $\EF$, the critical temperature is $T_c=0.567\,\Delta_F$, the ratio of weak coupling,
to within 0.002.

\emph{The law of growth.} In weak coupling the gap is proportional to $e^{-1/\lambda}$. With
$\lambda\propto e^{\td}$,
\begin{equation}
\frac{d\ln\Delta_F}{d\td}=\frac{1}{\lambda}.
\label{eq:growth}
\end{equation}
In the computed histories the product $\lambda\,d\ln\Delta_F/d\td$ lies between 1.00 and 1.03
wherever the gap is between $10^{-6}$ and $0.02$ of $\EF$.

Equation (\ref{eq:growth}) is the central quantitative statement about the superfluid on a deflating
spacetime. While the coupling grows by one e-folding, the gap of a weakly paired state with
$\lambda=0.1$ grows by ten e-foldings, a factor of twenty thousand. A paired state that is
unobservably weak at one time is strong a few e-foldings of the volume later.

\begin{figure*}
\includegraphics[width=\textwidth]{fig_gap_history}
\caption{\label{fig:gap}The gap along the deflation clock, for four sets of the mass of the radial
mode and of the Fermi momentum (in units of $m$). (a) The gap of the vector channel in units of the
Fermi kinetic energy. (b) The ratio of the critical temperature to the gap. (c) The rate of growth of
the gap against $1/\lambda$; the line is Eq.~(\ref{eq:growth}).}
\end{figure*}

